\subsection{The LEGB Rule Invariant}

When Python sees an unqualified name such as \texttt{x}, \texttt{print}, \texttt{len}, or \texttt{answer}, it must decide which object that name refers to. The search follows a stable order known as LEGB:

\begin{verbatim}
L = Local
E = Enclosing
G = Global
B = Built-in
\end{verbatim}

This rule applies to ordinary name lookup. It does not mean Python searches every object in memory. It searches a chain of namespaces connected to the current execution point.

The four levels are:

\begin{itemize}
    \item \textbf{Local}: names inside the currently executing function frame;
    \item \textbf{Enclosing}: names in outer function frames for nested functions;
    \item \textbf{Global}: names in the current module namespace;
    \item \textbf{Built-in}: names from the built-in namespace, such as \texttt{len}, \texttt{print}, \texttt{dict}, and \texttt{TypeError}.
\end{itemize}

The rule is not decorative terminology. It is the actual direction of lookup for unqualified names.

\subsubsection{Local (L): Active Function-Frame Bindings and CPython Fast-Local Storage}

The local namespace is the first place Python checks when a function body refers to a name.

\begin{verbatim}
def show_local():
    quote = "D'oh!"
    answer = 42

    print(f"quote:  {quote!r}")
    print(f"answer: {answer!r}")

show_local()
\end{verbatim}

Possible output:

\begin{verbatim}
quote:  "D'oh!"
answer: 42
\end{verbatim}

Inside \texttt{show\_local()}, the names \texttt{quote} and \texttt{answer} are local names. They belong to the active function call frame.

They are not global names.

\begin{verbatim}
def show_local():
    quote = "D'oh!"
    answer = 42

    print("inside function")
    print(f"'quote' in locals():  {'quote' in locals()}")
    print(f"'answer' in locals(): {'answer' in locals()}")

show_local()

print()
print("outside function")

try:
    print(quote)
except NameError as err:
    print(f"error type: {type(err)}")
    print(f"message:    {err}")
\end{verbatim}

Possible output:

\begin{verbatim}
inside function
'quote' in locals():  True
'answer' in locals(): True

outside function
error type: <class 'NameError'>
message:    name 'quote' is not defined
\end{verbatim}

The local names existed while the function frame was active. After the call ended, that frame was gone.

The \texttt{locals()} function displays the current local namespace.

\begin{verbatim}
def show_local_namespace():
    name = "Homer"
    score = 42

    local_map = locals()

    print(f"local_map:       {local_map}")
    print(f"type(local_map): {type(local_map)}")

show_local_namespace()
\end{verbatim}

Possible output:

\begin{verbatim}
local_map:       {'name': 'Homer', 'score': 42}
type(local_map): <class 'dict'>
\end{verbatim}

The displayed local namespace is dictionary-like.

\begin{verbatim}
def inspect_local_namespace():
    name = "Homer"
    score = 42

    local_map = locals()

    print(f"type(local_map): {type(local_map)}")
    print()

    selected_names = [
        "keys",
        "values",
        "items",
        "get",
        "__contains__",
        "__getitem__",
    ]

    for attr in selected_names:
        print(f"{attr:<14} {attr in dir(local_map)}")

    print()
    print(f"keys:              {list(local_map.keys())}")
    print(f"local_map['name']: {local_map['name']!r}")
    print(f"local_map.get('score'): {local_map.get('score')}")

inspect_local_namespace()
\end{verbatim}

Possible output:

\begin{verbatim}
type(local_map): <class 'dict'>

keys           True
values         True
items          True
get            True
__contains__   True
__getitem__    True

keys:              ['name', 'score']
local_map['name']: 'Homer'
local_map.get('score'): 42
\end{verbatim}

This display is useful for inspection, but it should not be confused with the exact internal storage strategy used by CPython for function locals.

In CPython, ordinary function local variables are stored in an optimized internal layout often called \textit{fast locals}. The function's compiled code object already knows the local variable names.

\begin{verbatim}
def sample_func(name, score):
    message = f"{name}: {score}"
    return message

code_obj = sample_func.__code__

print(f"co_varnames:   {code_obj.co_varnames}")
print(f"co_argcount:   {code_obj.co_argcount}")
print(f"co_nlocals:    {code_obj.co_nlocals}")
\end{verbatim}

Possible output:

\begin{verbatim}
co_varnames:   ('name', 'score', 'message')
co_argcount:   2
co_nlocals:    3
\end{verbatim}

The code object says that the function frame needs local slots for:

\begin{verbatim}
name
score
message
\end{verbatim}

That is why local variable access inside a function is not just a normal dictionary lookup every time. CPython can use the compiled local-variable layout.

We can inspect the code object itself.

\begin{verbatim}
def sample_func(name, score):
    message = f"{name}: {score}"
    return message

code_obj = sample_func.__code__

print(f"type(code_obj): {type(code_obj)}")
print()

selected_names = [
    "co_varnames",
    "co_argcount",
    "co_nlocals",
    "co_consts",
    "co_names",
    "co_code",
]

for attr in selected_names:
    print(f"{attr:<14} {attr in dir(code_obj)}")
\end{verbatim}

Possible output:

\begin{verbatim}
type(code_obj): <class 'code'>

co_varnames    True
co_argcount    True
co_nlocals     True
co_consts      True
co_names       True
co_code        True
\end{verbatim}

The local namespace therefore has two useful views:

\begin{itemize}
    \item the programmer-facing view: local names visible through \texttt{locals()};
    \item the CPython execution view: optimized local slots described by the code object.
\end{itemize}

For lookup behavior, the important rule is simpler:

\begin{quote}
If a name exists in the active local function frame, Python uses that local binding before looking anywhere else.
\end{quote}

This can shadow a global name.

\begin{verbatim}
answer = 100

def show_answer():
    answer = 42
    print(f"inside function: answer={answer}")

show_answer()

print(f"outside function: answer={answer}")
\end{verbatim}

Possible output:

\begin{verbatim}
inside function: answer=42
outside function: answer=100
\end{verbatim}

Inside the function, \texttt{answer} is found locally first. Python does not continue outward to the global \texttt{answer} because the local one already satisfies the lookup.

\subsubsection{Enclosing (E): Looking Upwards Through Cell Variables of Nested Lexical Scopes}

The enclosing level exists when a function is defined inside another function.

\begin{verbatim}
def outer():
    quote = "Nobody expects the Spanish Inquisition!"

    def inner():
        print(f"inner sees quote: {quote}")

    inner()

outer()
\end{verbatim}

Possible output:

\begin{verbatim}
inner sees quote: Nobody expects the Spanish Inquisition!
\end{verbatim}

The name \texttt{quote} is not local to \texttt{inner()}. It is local to \texttt{outer()}. Since \texttt{inner()} is nested inside \texttt{outer()}, Python can search the enclosing function scope.

The lookup path inside \texttt{inner()} is:

\begin{verbatim}
1. Local scope of inner()
2. Enclosing scope of outer()
3. Global module scope
4. Built-in scope
\end{verbatim}

This is lexical nesting. It depends on where the function is defined in the source, not merely where it is called.

\begin{verbatim}
quote = "global quote"

def outer():
    quote = "outer quote"

    def inner():
        print(f"quote: {quote}")

    inner()

outer()
\end{verbatim}

Possible output:

\begin{verbatim}
quote: outer quote
\end{verbatim}

The name \texttt{quote} is not found inside \texttt{inner()}. Python then checks the enclosing function scope and finds \texttt{outer}'s \texttt{quote}. It does not use the global \texttt{quote} because the enclosing binding was found first.

If the inner function has its own local name, that local name wins.

\begin{verbatim}
quote = "global quote"

def outer():
    quote = "outer quote"

    def inner():
        quote = "inner quote"
        print(f"quote: {quote}")

    inner()

outer()
\end{verbatim}

Possible output:

\begin{verbatim}
quote: inner quote
\end{verbatim}

The local level is still first.

A nested function that refers to an outer variable causes CPython to store that outer variable in a cell so the inner function can access it.

\begin{verbatim}
def outer():
    quote = "It is only a model."

    def inner():
        return quote

    return inner

fn = outer()

print(f"fn(): {fn()}")
print(f"fn.__closure__: {fn.__closure__}")
print(f"fn.__code__.co_freevars: {fn.__code__.co_freevars}")
\end{verbatim}

Possible output:

\begin{verbatim}
fn(): It is only a model.
fn.__closure__: (<cell at 0x...: str object at 0x...>,)
fn.__code__.co_freevars: ('quote',)
\end{verbatim}

The outer function has already finished, but the returned inner function still remembers the needed enclosing value. That preserved value lives through a closure cell.

The cell object can be inspected.

\begin{verbatim}
def outer():
    quote = "It is only a model."

    def inner():
        return quote

    return inner

fn = outer()
cell = fn.__closure__[0]

print(f"cell:       {cell}")
print(f"type(cell): {type(cell)}")
print()

selected_names = [
    "cell_contents",
    "__class__",
    "__eq__",
    "__repr__",
]

for attr in selected_names:
    print(f"{attr:<16} {attr in dir(cell)}")

print()
print(f"cell.cell_contents: {cell.cell_contents!r}")
print(f"type(contents):     {type(cell.cell_contents)}")
\end{verbatim}

Possible output:

\begin{verbatim}
cell:       <cell at 0x...: str object at 0x...>
type(cell): <class 'cell'>

cell_contents    True
__class__        True
__eq__           True
__repr__         True

cell.cell_contents: 'It is only a model.'
type(contents):     <class 'str'>
\end{verbatim}

This is the mechanical basis of enclosing-scope lookup. The inner function does not search an arbitrary dead stack frame. CPython preserves the required enclosing references in closure cells.

The outer function's code object also records which local names become cell variables.

\begin{verbatim}
def outer():
    quote = "It is only a model."

    def inner():
        return quote

    return inner

print(f"outer co_cellvars: {outer.__code__.co_cellvars}")

fn = outer()
print(f"inner co_freevars: {fn.__code__.co_freevars}")
\end{verbatim}

Possible output:

\begin{verbatim}
outer co_cellvars: ('quote',)
inner co_freevars: ('quote',)
\end{verbatim}

The name \texttt{quote} is a cell variable from the outer function's point of view. It is a free variable from the inner function's point of view.

The full closure mechanism will be studied later in the section on first-class functions and closures. Here, the important lookup rule is:

\begin{quote}
If a name is not local, Python looks through enclosing lexical function scopes before checking the module namespace.
\end{quote}

\subsubsection{Global (G): Module-Level Dictionary Namespaces and Active Script Execution State (\texttt{globals()})}

The global namespace is the namespace of the current module. In a simple script, it is the top-level script namespace.

A name assigned at top level is global to that module.

\begin{verbatim}
answer = 42
quote = "No soup for you!"

def show_globals():
    print(f"answer: {answer}")
    print(f"quote:  {quote}")

show_globals()
\end{verbatim}

Possible output:

\begin{verbatim}
answer: 42
quote:  No soup for you!
\end{verbatim}

Inside \texttt{show\_globals()}, the names \texttt{answer} and \texttt{quote} are not local. There is no enclosing function scope. Python therefore looks in the global module namespace.

The global namespace is visible through \texttt{globals()}.

\begin{verbatim}
answer = 42
quote = "No soup for you!"

global_map = globals()

print(f"type(global_map): {type(global_map)}")
print(f"'answer' in globals(): {'answer' in global_map}")
print(f"'quote' in globals():  {'quote' in global_map}")
print()
print(f"global_map['answer']: {global_map['answer']}")
print(f"global_map['quote']:  {global_map['quote']!r}")
\end{verbatim}

Possible output:

\begin{verbatim}
type(global_map): <class 'dict'>
'answer' in globals(): True
'quote' in globals():  True

global_map['answer']: 42
global_map['quote']:  'No soup for you!'
\end{verbatim}

The object returned by \texttt{globals()} is a dictionary.

\begin{verbatim}
answer = 42

global_map = globals()

print(f"type(global_map): {type(global_map)}")
print()

selected_names = [
    "keys",
    "values",
    "items",
    "get",
    "__contains__",
    "__getitem__",
    "__setitem__",
]

for attr in selected_names:
    print(f"{attr:<14} {attr in dir(global_map)}")

print()
print(f"global_map.get('answer'): {global_map.get('answer')}")
\end{verbatim}

Possible output:

\begin{verbatim}
type(global_map): <class 'dict'>

keys           True
values         True
items          True
get            True
__contains__   True
__getitem__    True
__setitem__    True

global_map.get('answer'): 42
\end{verbatim}

This dictionary is the module's namespace dictionary. Top-level assignments place names into it.

\begin{verbatim}
answer = 42

print(f"before: globals()['answer'] = {globals()['answer']}")

globals()["answer"] = 100

print(f"after:  answer = {answer}")
\end{verbatim}

Possible output:

\begin{verbatim}
before: globals()['answer'] = 42
after:  answer = 100
\end{verbatim}

This works because \texttt{globals()} returns the actual module namespace dictionary.

That does not mean ordinary code should casually mutate \texttt{globals()}. It is usually clearer to use normal assignment at module level. The point here is architectural: the global namespace is a dictionary-like object used for module-level name resolution.

A function can read a global name without any special declaration.

\begin{verbatim}
answer = 42

def show_answer():
    print(f"answer inside function: {answer}")

show_answer()
\end{verbatim}

Possible output:

\begin{verbatim}
answer inside function: 42
\end{verbatim}

But assignment inside a function creates a local binding by default. It does not automatically mutate the global namespace.

\begin{verbatim}
answer = 42

def change_answer():
    answer = 100
    print(f"inside function: answer={answer}")

change_answer()

print(f"outside function: answer={answer}")
\end{verbatim}

Possible output:

\begin{verbatim}
inside function: answer=100
outside function: answer=42
\end{verbatim}

The assignment inside \texttt{change\_answer()} created a local \texttt{answer}. The global \texttt{answer} remained unchanged.

The \texttt{global} declaration, which allows assignment to target the module namespace, belongs to the next subsection. For now, the lookup rule is:

\begin{quote}
If a name is not local and not found in an enclosing function scope, Python looks in the current module's global namespace.
\end{quote}

\subsubsection{Built-in (B): The Outer Built-In Namespace Boundary and the \texttt{builtins} Module}

If Python does not find a name in the local, enclosing, or global namespace, it checks the built-in namespace.

This is why names such as \texttt{print}, \texttt{len}, \texttt{dict}, \texttt{list}, \texttt{int}, and \texttt{TypeError} are normally available without import statements.

\begin{verbatim}
text = "D'oh!"

print(f"len(text): {len(text)}")
print(f"type(text): {type(text)}")
\end{verbatim}

Possible output:

\begin{verbatim}
len(text): 5
type(text): <class 'str'>
\end{verbatim}

The names \texttt{print}, \texttt{len}, and \texttt{type} are not defined locally by this code. They are found in the built-in namespace.

The built-in namespace is exposed through the standard \texttt{builtins} module.

\begin{verbatim}
import builtins

print(f"builtins.len:   {builtins.len}")
print(f"builtins.print: {builtins.print}")
print(f"builtins.dict:  {builtins.dict}")
print()
print(f"type(builtins): {type(builtins)}")
\end{verbatim}

Possible output:

\begin{verbatim}
builtins.len:   <built-in function len>
builtins.print: <built-in function print>
builtins.dict:  <class 'dict'>

type(builtins): <class 'module'>
\end{verbatim}

The module object can be inspected.

\begin{verbatim}
import builtins

print(f"type(builtins): {type(builtins)}")
print()

selected_names = [
    "len",
    "print",
    "dict",
    "list",
    "int",
    "TypeError",
    "__dict__",
]

for attr in selected_names:
    print(f"{attr:<12} {attr in dir(builtins)}")
\end{verbatim}

Possible output:

\begin{verbatim}
type(builtins): <class 'module'>

len          True
print        True
dict         True
list         True
int          True
TypeError    True
__dict__     True
\end{verbatim}

A module has a namespace dictionary.

\begin{verbatim}
import builtins

builtins_map = builtins.__dict__

print(f"type(builtins_map): {type(builtins_map)}")
print()
print(f"'len' in builtins_map:   {'len' in builtins_map}")
print(f"'print' in builtins_map: {'print' in builtins_map}")
print(f"'dict' in builtins_map:  {'dict' in builtins_map}")
print()
print(f"builtins_map['len']: {builtins_map['len']}")
\end{verbatim}

Possible output:

\begin{verbatim}
type(builtins_map): <class 'dict'>

'len' in builtins_map:   True
'print' in builtins_map: True
'dict' in builtins_map:  True

builtins_map['len']: <built-in function len>
\end{verbatim}

The built-in function object \texttt{len} can also be inspected.

\begin{verbatim}
import builtins

fn = builtins.len

print(f"fn:       {fn}")
print(f"type(fn): {type(fn)}")
print()

selected_names = [
    "__call__",
    "__name__",
    "__doc__",
    "__self__",
    "__text_signature__",
]

for attr in selected_names:
    print(f"{attr:<22} {attr in dir(fn)}")

print()
print(f"fn.__name__: {fn.__name__!r}")
print(f"fn('D\\'oh!'): {fn(\"D'oh!\")}")
\end{verbatim}

Possible output:

\begin{verbatim}
fn:       <built-in function len>
type(fn): <class 'builtin_function_or_method'>

__call__                True
__name__                True
__doc__                 True
__self__                True
__text_signature__      True

fn.__name__: 'len'
fn('D\'oh!'): 5
\end{verbatim}

Built-in names can be shadowed by local or global names. This is usually a bad idea, but it clearly shows LEGB order.

\begin{verbatim}
len = 42

print(f"global len: {len}")
print(f"type(len):  {type(len)}")
print()

try:
    print(len("D'oh!"))
except TypeError as err:
    print(f"error type: {type(err)}")
    print(f"message:    {err}")
\end{verbatim}

Possible output:

\begin{verbatim}
global len: 42
type(len):  <class 'int'>

error type: <class 'TypeError'>
message:    'int' object is not callable
\end{verbatim}

The global name \texttt{len} now points to the integer object \texttt{42}. When Python sees \texttt{len("D'oh!")}, it finds the global \texttt{len} before it reaches the built-in namespace. It then tries to call the integer object, which fails.

The original built-in is still available through \texttt{builtins.len}.

\begin{verbatim}
import builtins

len = 42

print(f"global len:       {len}")
print(f"builtins.len(...): {builtins.len(\"D'oh!\")}")
\end{verbatim}

Possible output:

\begin{verbatim}
global len:       42
builtins.len(...): 5
\end{verbatim}

A local name can also shadow a built-in name.

\begin{verbatim}
def show_shadowing():
    list = "This is not the list class."

    print(f"local list:       {list!r}")
    print(f"type(local list): {type(list)}")

    try:
        values = list((1, 2, 3))
    except TypeError as err:
        print(f"error type: {type(err)}")
        print(f"message:    {err}")

show_shadowing()
\end{verbatim}

Possible output:

\begin{verbatim}
local list:       'This is not the list class.'
type(local list): <class 'str'>
error type: <class 'TypeError'>
message:    'str' object is not callable
\end{verbatim}

Inside the function, the local name \texttt{list} hides the built-in class \texttt{list}. Python finds the local string first and stops searching.

The compact rule is:

\begin{quote}
For an unqualified name, Python searches Local, then Enclosing, then Global, then Built-in. The first binding found is the one used.
\end{quote}

A complete example can show all four levels.

\begin{verbatim}
import builtins

name = "global name"

def outer():
    name = "enclosing name"

    def inner():
        name = "local name"

        print(f"name:          {name}")
        print(f"builtins.len:  {builtins.len('D\\'oh!')}")

    inner()

outer()
\end{verbatim}

Possible output:

\begin{verbatim}
name:          local name
builtins.len:  5
\end{verbatim}

If the local \texttt{name} is removed, the enclosing one is found. If the enclosing one is removed, the global one is found. If the global one is removed and no built-in exists with that name, lookup fails with \texttt{NameError}.

That is the LEGB invariant.